\section*{Capítulo \ref{ch:b3:probability} --- Probabilidad:
fundamentos y ley de los grandes números}

\begin{solution}{exo:b3:probability:1}
(a) Para $u \in \intoo01$: $\P(F(X) \leq u) = \P(X \leq
F^{-1}(u)) = F(F^{-1}(u)) = u$ (la
\omterm{def:b3:topology:continuity}{continuidad} y la monotonía estricta
hacen de $F$ una biyección sobre $\intoo01$ con
$\{F(X) \leq u\} = \{X \leq F^{-1}(u)\}$): $F(X)$ es uniforme.
Recíprocamente, $\P(G(U) \leq t) = \P(U \leq F(t)) = F(t)$: para simular
una \omterm{def:b3:probability:space}{ley}, aplíquese la inversa de la
función de distribución a una muestra uniforme.

(b) $Y = X^2$, $X$ uniforme en $\intcc{-1}1$: para $t \in
\intcc01$, $F_Y(t) = \P(-\sqrt t \leq X \leq \sqrt t) = \sqrt
t$: \omterm{ex:b3:lebesgue:gamma}{densidad}
$\frac1{2\sqrt t}\mathbf 1_{\intoo01}$. Y
$\P\bigl(-\frac1\lambda\ln U \leq t\bigr) = \P(U \geq
\eu^{-\lambda t}) = 1 - \eu^{-\lambda t}$: la exponencial
$\mathcal E(\lambda)$ --- la inversión en acción.
\end{solution}

\begin{solution}{exo:b3:probability:2}
(a) Poisson: $\E X = \sum_{k\geq0}k\,\eu^{-\lambda}
\frac{\lambda^k}{k!} = \lambda$, $\E[X(X-1)] = \lambda^2$, de modo que
$\V = \lambda^2 + \lambda - \lambda^2 = \lambda$. Geométrica
($\P(X = k) = p(1-p)^{k-1}$): $\E X = \frac1p$, $\V =
\frac{1-p}{p^2}$ (derívese dos veces la serie geométrica).

(b) $G(t) = \P(T > t)$ es no creciente con $G(0^+)\dots$
$G \colon \intco0\infty \to \intoc01$; la falta de memoria se lee
$G(t + s) = G(t)G(s)$. Entonces $G(n t) = G(t)^n$ y $G(t/n) =
G(t)^{1/n}$: $G(q) = G(1)^q$ para $q \geq 0$ racional; escribiendo
$G(1) = \eu^{-\lambda}$ ($\in \intoo01$: $G(1) = 1$ forzaría
$G \equiv 1$, imposible para una
\omterm{def:b3:probability:space}{variable aleatoria} finita; $G(1) = 0$
queda excluido por hipótesis) y encajonando un $t$ arbitrario entre
racionales (monotonía): $G(t) =
\eu^{-\lambda t}$ --- la \omterm{def:b3:probability:space}{ley}
exponencial. El recíproco es un cálculo.
\end{solution}

\begin{solution}{exo:b3:probability:3}
(a) $\sigma(f_i(X_i)) = f_i(X_i)^{-1}(\mathcal B) \subseteq
X_i^{-1}(\mathcal B) = \sigma(X_i)$ ($f_i$ borelianas), y las
sub-$\sigma$-álgebras de $\sigma$-álgebras
\omterm{def:b3:probability:independence}{independientes} son
\omterm{def:b3:probability:independence}{independientes} (la identidad
que las define vale a fortiori).

(b) $\sigma(A_i) = \{\varnothing, A_i, A_i^c, \Omega\} =
\sigma(A_i^c) = \sigma(\mathbf 1_{A_i})$: las tres afirmaciones expresan
la \omterm{def:b3:probability:independence}{independencia} de las mismas
$\sigma$-álgebras. (Que la factorización sobre los $A_i$ se propague a
los complementarios es el argumento del $\lambda$-sistema contenido en
la equivalencia de la~\cref{def:b3:probability:independence} --- o bien
la inclusión-exclusión directa.)

(c) $\P(A) = \P(B) = \P(C) = \frac12$; $A\cap B = A\cap C =
B\cap C$ sobre los pares: cada intersección es «ambas caras» o análoga,
de probabilidad $\frac14$:
\omterm{def:b3:probability:independence}{independientes} dos a dos. Pero
$\P(A\cap B\cap C) = \P(\text{CC}) = \frac14 \neq
\frac18$: no \omterm{def:b3:probability:independence}{independientes}
--- $C$ queda determinado por $A$ y
$B$.
\end{solution}

\begin{solution}{exo:b3:probability:4}
(a) Sean $T$ el texto, de longitud $L$, y $q = a^{-L}$ ($a$ el tamaño
del alfabeto). Los sucesos $E_k = \{$las posiciones $kL+1,
\dots, (k+1)L$ deletrean $T\}$ son
\omterm{def:b3:probability:independence}{independientes} (bloques
disjuntos de letras i.i.d.), cada uno de probabilidad $q > 0$:
$\sum\P(E_k)
= \infty$, y Borel--Cantelli (2) da infinitas apariciones c.s.

(b) Cota superior: $\P\bigl(R_n \geq (1+\varepsilon)\log_2n\bigr)
\leq 2^{-(1+\varepsilon)\log_2n} = n^{-(1+\varepsilon)}$, sumable: por
Borel--Cantelli (1), c.s.\ solo hay un número finito de tales $n$. Cota
inferior: empaquétense bloques disjuntos --- el $j$-ésimo, de longitud
$\ell_j = \lceil\log_2s_j\rceil$, empezando en $s_j =
\sum_{i<j}\ell_i$; los sucesos «el bloque $j$ es todo unos» son
\omterm{def:b3:probability:independence}{independientes} de probabilidad
$2^{-\ell_j} \asymp
\frac1{s_j} \asymp \frac1{j\log_2 j}$, cuya suma diverge:
Borel--Cantelli (2) da infinitos bloques de unos, es decir,
$R_{s_j} \geq \log_2 s_j$ infinitas veces. En conjunto: la longitud
máxima de racha entre los $n$ primeros dígitos es
$(1 + o(1))\log_2n$ a.s.
\end{solution}

\begin{solution}{exo:b3:probability:5}
(a) $R = \bigl(\limsup\abs{X_n}^{1/n}\bigr)^{-1}$ no cambia si se
modifican un número finito de $X_n$: para todo $N$, $R$ es
$\sigma(X_N, X_{N+1}, \dots)$-medible, es decir,
\omterm{def:b3:lebesgue:measurable}{medible} respecto de la cola.
Entonces cada suceso $\{R \leq c\}$ tiene probabilidad $0$ o $1$
(el~\cref{thm:b3:probability:zeroone}), de modo que la función de
distribución de $R$ solo toma los valores $0, 1$: salta en un único
punto $c_0 \in
\intcc0{+\infty}$, y $R = c_0$ c.s.

(b) La convergencia de $\sum X_n$ y la de $\frac{S_n}n$ son insensibles
a cambiar un número finito de términos (para la segunda: los términos
modificados aportan $O(1/n) \to 0$): sucesos de cola;
\omterm{thm:b3:probability:zeroone}{ley cero--uno}.

(c) $\{X_1 > 0\}$ depende de $X_1$: para signos i.i.d.\
($\P(X_1 = \pm1) = \frac12$), su probabilidad vale $\frac12
\notin \{0,1\}$ --- ninguna contradicción, no es un suceso de cola.
\end{solution}

\begin{solution}{exo:b3:probability:6}
Trabájese en $(\intcc01, \lambda)$. (a) La máquina de escribir
$\mathbf 1_{I_n}$
(\cref{exo:b3:lp:3}): $\norm{X_n}_p^p = \lambda(I_n) \to 0$
(en todo $p < \infty$), luego también en probabilidad; en cada $\omega$
los valores $0$ y $1$ reaparecen ambos: no hay convergencia puntual en
ningún punto. (b) $X_n = n\mathbf 1_{\intoo0{1/n}} \to 0$ fuera de $0$,
pero
$\norm{X_n}_p \geq n^{1 - 1/p} \geq 1$.
(c) $X_n = \sqrt n\,\mathbf 1_{\intoo0{1/n}}$: $\E\abs{X_n} =
n^{-1/2} \to 0$, $\E X_n^2 = 1$.
(d) De cualquier subsucesión extráigase (convergencia en probabilidad)
una subsucesión ulterior que converja c.s.\
(la~\cref{prop:b3:probability:modes}(c)); la convergencia dominada da
convergencia en $L^1$ a lo largo de ella, con el \emph{mismo} límite
$X$. Así, toda subsucesión de la sucesión numérica $\E\abs{X_n - X}$
tiene una subsubsucesión que tiende a $0$: la sucesión entera tiende a
$0$.
\end{solution}

\begin{solution}{exo:b3:probability:7}
(a) $\hat p_n = \frac{S_n}n$ con $S_n$ binomial: $\V(\hat
p_n) = \frac{p(1-p)}n \leq \frac1{4n}$, y Chebyshev
(la~\cref{prop:b3:probability:markov}) da la cota. (b) Resuélvase
$\frac1{4n(0.03)^2} \leq 0.05$: $n \geq
\frac{1}{4\cdot0.0009\cdot0.05} \approx 5556$. El teorema central del
límite justificará $n \approx 1070$ para la misma garantía: Chebyshev
paga su generalidad con un factor
$\approx 5$.
\end{solution}

\begin{solution}{exo:b3:probability:8}
(a) Si $S_n \sim \mathcal B(n, x)$, el teorema de transferencia da
$\E\bigl[f(\frac{S_n}n)\bigr] =
\sum_k\binom nkx^k(1-x)^{n-k}f(\frac kn) = B_nf(x)$.

(b)--(c) Sea $\omega = \omega_f$ el módulo de
\omterm{def:b3:topology:continuity}{continuidad}
($\abs{f(u) - f(v)} \leq \omega(\abs{u - v})$, y $\omega(c
\delta) \leq (1 + c)\,\omega(\delta)$ encadenando pasos). Entonces, para
todo $\delta > 0$,
\[
\abs{f(u) - f(x)} \leq \Bigl(1 + \frac{(u -
x)^2}{\delta^2}\Bigr)\omega(\delta)
\]
(si $\abs{u - x} \leq \delta$, es claro; en caso contrario,
$\omega(\abs{u-x}) \leq (1 + \frac{\abs{u-x}}\delta)
\omega(\delta) \leq (1 + \frac{(u-x)^2}{\delta^2})
\omega(\delta)$). Tómense \omterm{def:b3:probability:space}{esperanzas}
en $u = \frac{S_n}n$:
\[
\abs{B_nf(x) - f(x)} \leq
\Bigl(1 + \frac{\V(S_n/n)}{\delta^2}\Bigr)\omega(\delta)
\leq \Bigl(1 + \frac{1}{4n\delta^2}\Bigr)\omega(\delta) ;
\]
con $\delta = n^{-1/2}$: $\norm{B_nf - f}_\infty \leq
\frac54\,\omega\bigl(n^{-1/2}\bigr) \leq
\frac32\,\omega\bigl(n^{-1/2}\bigr) \to 0$
(\omterm{def:b3:topology:continuity}{continuidad} uniforme en el
\omterm{def:b3:topology:compact}{compacto}): un teorema de Weierstrass
probabilístico, con velocidad explícita y uniforme.
\end{solution}

\begin{solution}{exo:b3:probability:9}
(a) Una vez recogidos $k - 1$ tipos, cada extracción es nueva con
probabilidad $p_k = \frac{n-k+1}n$: $\tau_k$ es geométrica $(p_k)$, y
las $\tau_k$ son
\omterm{def:b3:probability:independence}{independientes} (las
extracciones lo son). Sumas:
$\E T_n = \sum_k\frac n{n-k+1} = nH_n \sim n\ln n$;
$\V(T_n) = \sum\frac{1 - p_k}{p_k^2} \leq
n^2\sum_{j=1}^n\frac1{j^2} \leq \frac{\pi^2}6n^2$.

(b) Chebyshev: $\P\bigl(\abs{T_n - nH_n} \geq \varepsilon
n\ln n\bigr) \leq \frac{\pi^2n^2/6}{\varepsilon^2n^2\ln^2n}
\to 0$, y $\frac{nH_n}{n\ln n} \to 1$: $\frac{T_n}{n\ln n}
\to 1$ en probabilidad.

(c) Cota de la unión: $T_n > t$ significa que algún tipo no ha salido
tras $\lceil t\rceil$ extracciones, de modo que $\P(T_n > t) \leq n(1 -
\frac1n)^{t} \leq n\,\eu^{-t/n}$; en $t = \beta n\ln n$:
$\leq n^{1 - \beta}$. Para $\beta > 1$, $\sum_m
2^{m(1-\beta)} < \infty$: Borel--Cantelli da, a lo largo de $n =
2^m$, que c.s.\ $T_n \leq \beta n\ln n$ a partir de cierto índice --- en
particular, para todo $\beta > 2$ como se enunció (cualquier $\beta > 1$
sirve a lo largo de la subsucesión).
\end{solution}

\begin{solution}{exo:b3:probability:10}
(a) Los dígitos de índice par $(b_{2k})_k$ son bits equilibrados i.i.d.\
(una subfamilia de la familia
independiente de dígitos), de
modo que $U =
\sum_kb_{2k}2^{-k}$ da a cada intervalo diádico su probabilidad correcta
(como en el~\cref{thm:b3:probability:existence}): uniforme; análogamente
$V$; y $(U, V)$ dependen de bloques de dígitos disjuntos:
\omterm{def:b3:probability:independence}{independientes} (factorización
en rectángulos diádicos y después Dynkin).

(b) $\Phi(\omega) = (U(\omega), V(\omega))$ es
\omterm{def:b3:lebesgue:measurable}{medible} con
$\Phi_*\lambda = \lambda\otimes\lambda = \lambda_2$ (coincidencia en los
rectángulos diádicos más unicidad). Entrelazar dígitos define una
inversa definida fuera del conjunto (nulo) de los racionales diádicos de
cualquiera de los dos factores: una biyección que conserva la \omterm{def:b3:measure:measure}{medida}
entre subconjuntos de \omterm{def:b3:measure:measure}{medida} total de
$\intcc01$ y $\intcc01^2$. Contrástese con Peano
(el~\cref{pb:b3:topology:1}): la
\omterm{def:b3:topology:continuity}{continuidad} forzaba la
sobreyectividad sin inyectividad; renunciar a la
\omterm{def:b3:topology:continuity}{continuidad} a cambio de la mera
\omterm{def:b3:lebesgue:measurable}{medibilidad} compra un isomorfismo
de \omterm{def:b3:measure:measure}{medida} --- la dimensión es invisible para la teoría de la
\omterm{def:b3:measure:measure}{medida} y visible para la
\omterm{def:b3:topology:topology}{topología}.
\end{solution}

\begin{solution}{exo:b3:probability:11}
(a) La \omterm{def:b3:topology:continuity}{continuidad} de la
distribución hace nulos los sucesos de empate (como en los argumentos de
estadísticos de orden del capítulo), y las $n!$ ordenaciones relativas
de $(X_1, \dots, X_n)$ son intercambiables, luego igualmente probables.
$R_n = 1$ significa que el máximo ocupa la última posición: probabilidad
$\frac{(n-1)!}{n!} = \frac1n$.

(b) Fíjese $n$ y condiciónese al orden relativo de $X_1,
\dots, X_{n-1}$: insertar $X_n$ en una de las $n$ posiciones de rango
posibles es uniforme e
independiente de ese orden
(intercambiabilidad de la $n$-tupla). Por tanto, $R_n$ (el suceso «$X_n$
ocupa la primera posición») es
independiente de toda la
historia de récords $(R_1, \dots, R_{n-1})$, que es función del orden
relativo de las $n - 1$ primeras variables. La inducción da la
\omterm{def:b3:probability:independence}{independencia} \omterm{def:b3:complete:complete}{completa} con
$\P(R_n = 1) =
\frac1n$.

(c) $\sum\P(R_n = 1) = \sum\frac1n = \infty$ con
\omterm{def:b3:probability:independence}{independencia}: la segunda
mitad de Borel--Cantelli da infinitos récords c.s.\ (los récords nunca
cesan --- pero se ralean logarítmicamente:
$\E[\#\text{récords} \leq n] =
H_n \approx \ln n$). Récords consecutivos: $\P(R_n = R_{n+1}
= 1) = \frac1{n(n+1)}$
(\omterm{def:b3:probability:independence}{independencia}), y
\[
\sum_n\frac1{n(n+1)} = \sum_n\Bigl(\frac1n -
\frac1{n+1}\Bigr) = 1 < \infty :
\]
se aplica la primera mitad de Borel--Cantelli --- solo hay, c.s., un
número finito de pares de récords consecutivos.
\end{solution}

\begin{solution}{exo:b3:probability:12}
(a) Una racha de longitud $\ell$ que empieza en la posición $i \leq n$
tiene probabilidad $2^{-\ell}$; cota de la unión: $\P(L_n \geq \ell)
\leq n2^{-\ell}$. Con $\ell_n = (1 +
\varepsilon)\log_2n$: $\P(L_n \geq \ell_n) \leq
n^{-\varepsilon}$. A lo largo de $n = 2^k$:
$\sum_k2^{-k\varepsilon} < \infty$, de modo que c.s.\ $L_{2^k} <
(1+\varepsilon)k$ a partir de cierto índice (Borel--Cantelli); para $n$
general, tómese $2^{k-1} < n \leq 2^k$ y úsense la monotonía de $L_n$ y
$\log_22^{k-1} \leq \log_2n$: $L_n \leq L_{2^k}
< (1 + \varepsilon)k \leq (1 + \varepsilon)\frac{k}{k-1}
\log_2n$, y el factor extra se absorbe agrandando ligeramente
$\varepsilon$.

(b) Con $\ell = \lceil(1 - \varepsilon)\log_2n\rceil$ y
$m = \lfloor n/\ell\rfloor$ bloques disjuntos: los bloques son
\omterm{def:b3:probability:independence}{independientes}, cada uno todo
caras con probabilidad $2^{-\ell}
\geq n^{-(1-\varepsilon)}/2$, de modo que
\[
\P(L_n < \ell) \leq \bigl(1 - 2^{-\ell}\bigr)^{m}
\leq \exp\bigl(-m2^{-\ell}\bigr)
\leq \exp\Bigl(-c\,\frac{n^{\varepsilon}}{\log_2n}\Bigr)
\]
para cierta constante $c > 0$ y $n$ grande. Estas probabilidades son
sumables a lo largo de $n = 2^k$ (de hecho, en todo $n$):
Borel--Cantelli da c.s.\ $L_n \geq (1 -
\varepsilon)\log_2n$ a partir de cierto índice (la monotonía rellena los
huecos entre los $2^k$ como en (a), sin daño).

(c) Ambas cotas a lo largo de una sucesión $\varepsilon = \frac1j$,
intersecando una cantidad numerable de sucesos de
\omterm{def:b3:measure:measure}{medida} total:
$\frac{L_n}{\log_2n} \to 1$ c.s. Para $n = 10^6$: $\log_2n
\approx 19.9$ --- una racha de $\approx 20$ caras no es una anomalía
sospechosa, sino una certeza matemática, y su ausencia es indicio de un
humano fingiendo «azar» (rara vez se atreve nadie a escribir más de $5$
o $6$ caras seguidas).
\end{solution}

\begin{solution}{pb:b3:probability:1}
\textbf{1.} $X_n^{\pm}$ son funciones borelianas de $X_n$: siguen siendo
\omterm{def:b3:probability:independence}{independientes} dos a dos
(el~\cref{exo:b3:probability:3}(a)) e idénticamente distribuidas e
\omterm{def:b3:lebesgue:l1}{integrables}, con $\E X_1 = \E X_1^+ - \E
X_1^-$. Si el teorema vale para variables no negativas, aplíquese a
ambas mitades y réstese:
$\frac{S_n}n = \frac{S_n^+}n - \frac{S_n^-}n \to \E X_1^+ -
\E X_1^- = m$ a.s.

\textbf{2.} $\P(X_n \neq Y_n) = \P(X_n > n) = \P(X_1 > n)$
(\omterm{def:b3:probability:space}{leyes} idénticas), y
$\sum_n\P(X_1 > n) \leq \sum_n\P(X_1
\geq n) \leq \E X_1 < \infty$ (el~\cref{exo:b3:product:3}(a)).
Borel--Cantelli (1): c.s.\ $X_n = Y_n$ para todo $n$ grande, de modo que
$S_n - S_n^*$ es finalmente constante en $n$:
$\frac{S_n - S_n^*}n \to 0$ c.s., y las dos sumas normalizadas comparten
su comportamiento asintótico.

\textbf{3.} $X_1\mathbf 1_{X_1 \leq n} \nearrow X_1$: el teorema de
convergencia monótona da $\E Y_n \to m$; las medias de Cesàro de una
sucesión convergente convergen al mismo límite: $\frac{\E S_n^*}n =
\frac1n\sum_{k\leq n}\E Y_k \to m$. Basta, por tanto, demostrar
$\frac{S^*_n - \E S^*_n}{n} \to 0$ c.s.

\textbf{4.} $\V(Y_n) \leq \E Y_n^2 = \E[X_1^2\mathbf
1_{X_1\leq n}]$. Por Tonelli para series,
\[
\sum_n\frac{\E[X_1^2\mathbf 1_{X_1\leq n}]}{n^2}
= \E\Bigl[X_1^2\!\!\sum_{n \geq \max(X_1, 1)}\!\frac1{n^2}
\Bigr]
\leq \E\Bigl[X_1^2\cdot\frac{4}{\max(X_1,1)}\Bigr]
\leq 4\,\E[X_1] < \infty,
\]
usando $\sum_{n\geq x}n^{-2} \leq \frac4x$ para $x \geq 1$ (para
$x \geq 2$: $\leq \frac1{x-1} \leq \frac2x$; para $1
\leq x < 2$: $\leq \frac{\pi^2}6 \leq \frac4x$, pues $\frac4x > 2$), y
$X_1^2/\max(X_1, 1) \leq X_1$ en ambos casos $X_1 \gtrless 1$.

\textbf{5.} La \omterm{def:b3:probability:independence}{independencia}
dos a dos da $\E[(Y_i - \E
Y_i)(Y_j - \E Y_j)] = 0$ para $i \neq j$ (la fórmula del producto para
dos variables), de modo que las varianzas se suman: $\V(S^*_k) =
\sum_{n\leq k}\V(Y_n)$. Chebyshev en cada $k_j$ y sumando:
\[
\sum_j\P\Bigl(\abs{S^*_{k_j} - \E S^*_{k_j}} \geq
\varepsilon k_j\Bigr)
\leq \frac1{\varepsilon^2}\sum_j\frac1{k_j^2}\sum_{n\leq
k_j}\V(Y_n)
= \frac1{\varepsilon^2}\sum_n\V(Y_n)\!\!\sum_{j : k_j\geq
n}\!\frac1{k_j^2}
\]
(Tonelli para la serie doble de términos no negativos).

\textbf{6.} $k_j = \lfloor\alpha^j\rfloor \geq
\frac{\alpha^j}2$ (válido en cuanto $\alpha^j \geq 1$, es decir, para
todo $j \geq 0$: $\lfloor x\rfloor \geq \frac x2$ si $x \geq
1$). Por tanto,
\[
\sum_{j : k_j \geq n}\frac1{k_j^2}
\leq 4\sum_{j : \alpha^j \geq n}\alpha^{-2j}
\leq \frac{4}{1 - \alpha^{-2}}\cdot\frac1{n^2}
= \frac{C_\alpha}{n^2},
\]
(serie geométrica desde el primer $j$ con $\alpha^j \geq
n$). Combinando con las preguntas 4--5, la suma doble es finita;
Borel--Cantelli (1), aplicado para cada racional $\varepsilon$ e
intersecado, da $\frac{S^*_{k_j} - \E S^*_{k_j}}{k_j} \to 0$ c.s.\ y,
con la pregunta 3, $\frac{S^*_{k_j}}{k_j} \to m$ c.s.

\textbf{7.} $Y_n \geq 0$ hace $n \mapsto S^*_n$ no decreciente: para
$k_j \leq n \leq k_{j+1}$,
\[
\frac{S^*_{k_j}}{k_{j+1}} \leq \frac{S^*_n}{n} \leq
\frac{S^*_{k_{j+1}}}{k_j},
\]
que es el sándwich exhibido tras insertar $\frac{k_j}{k_{j+1}}$ y
$\frac{k_{j+1}}{k_j}$. Como $\frac{k_{j+1}}{k_j} \to \alpha$, la
pregunta 6 da, c.s.,
\[
\frac m\alpha \leq \liminf_n\frac{S^*_n}n \leq
\limsup_n\frac{S^*_n}n \leq \alpha m .
\]

\textbf{8.} Aplíquese la pregunta 7 para $\alpha = 1 + \frac1p$, $p
\in \N^*$: una cantidad numerable de sucesos c.s.; en su intersección,
haciendo $p \to \infty$: $\lim\frac{S^*_n}n =
m$ c.s. Con las preguntas 1--3, $\frac{S_n}n \to \E X_1$ c.s.: la
\omterm{def:b3:probability:space}{ley} fuerte de los grandes números,
bajo \omterm{def:b3:probability:independence}{independencia} dos a dos.

\textbf{9.} Las hipótesis de tipo \omterm{def:b3:probability:independence}{independencia} aparecieron tres veces:
(i) la aditividad de las varianzas (pregunta 5) --- basta la de dos a
dos; (ii) la idéntica distribución, en las sumas de truncamiento
(pregunta 2) y en el cálculo de la media (pregunta 3) --- sin
\omterm{def:b3:probability:independence}{independencia} alguna; (iii)
Borel--Cantelli (1) (preguntas 2 y 6) --- válido sin ninguna
\omterm{def:b3:probability:independence}{independencia}. La
\omterm{def:b3:probability:independence}{independencia} mutua \omterm{def:b3:complete:complete}{completa}
no se invocó nunca: la observación de Etemadi.

\textbf{10.} Fíjense una base $b$ y un dígito $r$. Los dígitos en base
$b$, $(d_k)$, de una $\omega$ uniforme son i.i.d.\ uniformes en
$\{0, \dots, b-1\}$ (cada valor del vector de dígitos ocupa un intervalo
de longitud $b^{-m}$: el argumento
del~\cref{thm:b3:probability:existence}, palabra por palabra). La
\omterm{def:b3:probability:space}{ley} fuerte aplicada a las variables
acotadas i.i.d.\ $\mathbf
1_{d_k = r}$ da: c.s., la frecuencia del dígito $r$ tiende a $\frac1b$.
Intersecando en la cantidad numerable de pares $(b, r)$: casi todo
número es \emph{simplemente normal en toda base}. Un número explícito no
normal: $x = 0.100100100\ldots_2$ (frecuencia de unos
$\frac13 \neq \frac12$). El contraste es humillante: casi todos los
números son normales y, sin embargo, para $\sqrt2$, $\eu$ o $\pi$ la
normalidad sigue sin demostrarse --- la teoría de la
\omterm{def:b3:measure:measure}{medida} cuenta sin exhibir.

\textbf{11.} Por el~\cref{exo:b3:probability:10} iterado, una sola
variable uniforme produce una sucesión de \emph{vectores} uniformes
i.i.d.\ $U_k$ en $\intcc01^d$ (divídase el conjunto de dígitos de cada
$U_n$ del~\cref{thm:b3:probability:existence} en $d$ subfamilias). Para
$g \in L^1(\intcc01^d)$, las variables $g(U_k)$ son i.i.d.\
\omterm{def:b3:lebesgue:l1}{integrables} de media $\int g\,\dd\lambda_d$
(transferencia): la \omterm{def:b3:probability:space}{ley} fuerte da
\[
\frac1n\sum_{k=1}^ng(U_k)
\xrightarrow[n\to\infty]{\text{c.s.}}
\int_{\intcc01^d}g\,\dd\lambda_d :
\]
la integración de \omterm{ex:b3:probability:sllnapps}{Monte Carlo}
converge casi seguramente, en toda dimensión --- el tamaño del error es
asunto del teorema central del límite (el~\cref{ch:b3:clt}).

\textbf{12.} Sea $A_k = \{\abs{S_k} \geq \varepsilon\} \cap
\bigcap_{j<k}\{\abs{S_j} < \varepsilon\}$: los $A_k$ son disjuntos con
unión $A = \{\max_{k\leq n}\abs{S_k} \geq
\varepsilon\}$. Entonces
\[
\E S_n^2 \geq \sum_{k=1}^n\E\bigl[S_n^2\mathbf 1_{A_k}\bigr]
= \sum_{k=1}^n\E\Bigl[\bigl(S_k^2 + 2S_k(S_n - S_k) + (S_n
- S_k)^2\bigr)\mathbf 1_{A_k}\Bigr]
\geq \sum_{k=1}^n\E\bigl[S_k^2\mathbf 1_{A_k}\bigr],
\]
porque el término cruzado se anula: $S_k\mathbf 1_{A_k}$ es una función
boreliana de la coalición $(Z_1, \dots, Z_k)$, que es
independiente de $S_n - S_k$,
función de $(Z_{k+1},
\dots, Z_n)$ (el~\cref{thm:b3:probability:independence}), de modo que
$\E[S_k\mathbf 1_{A_k}(S_n - S_k)] = \E[S_k\mathbf
1_{A_k}]\,\E[S_n - S_k] = 0$. En $A_k$, $S_k^2 \geq
\varepsilon^2$, de donde $\E S_n^2 \geq
\varepsilon^2\sum_k\P(A_k) = \varepsilon^2\P(A)$; y $\E
S_n^2 = \sum_{k\leq n}\V(Z_k)$ (las varianzas se suman). El paso
decisivo es la factorización: $S_k\mathbf 1_{A_k}$ es una función
\emph{no lineal} de todo el primer bloque, y su
\omterm{def:b3:probability:independence}{independencia} del segundo
bloque es \omterm{def:b3:probability:independence}{independencia} de
coaliciones --- la
\omterm{def:b3:probability:independence}{independencia} dos a dos de las
$Z_i$ solo descorrelaciona pares y no la justificaría.

\textbf{13.} Fíjese $N$ y aplíquese la pregunta 12 a $Z_{N+1},
\dots, Z_{N+m}$:
\[
\P\Bigl(\max_{N < k \leq N+m}\abs{S_k - S_N} >
\varepsilon\Bigr) \leq
\frac1{\varepsilon^2}\sum_{j=N+1}^{N+m}\V(Z_j) \leq
\frac{r_N}{\varepsilon^2},
\qquad r_N = \sum_{j>N}\V(Z_j) .
\]
Los sucesos crecen con $m$; la
\omterm{def:b3:topology:continuity}{continuidad} por abajo da
$\P(\sup_{k>N}\abs{S_k - S_N} > \varepsilon) \leq
r_N/\varepsilon^2$, y $r_N \to 0$ por hipótesis. Por tanto, para cada
$p \in \N^*$, $\P\bigl(\bigcap_N\{\sup_{k>N}
\abs{S_k - S_N} > \frac1p\}\bigr) \leq \inf_Np^2r_N = 0$: casi
seguramente, para todo $p$ hay un $N$ con
$\sup_{k>N}\abs{S_k - S_N} \leq \frac1p$ (interséquense los sucesos
c.s.\ en $p$, en cantidad numerable), de modo que $\abs{S_k -
S_l} \leq \frac2p$ para todos $k, l > N$: las sumas parciales son c.s.\
de Cauchy, luego c.s.\ convergentes.

\textbf{14.} Las variables $Z_n = x_n\varepsilon_n$ son
\omterm{def:b3:probability:independence}{independientes} (funciones
borelianas de variables
\omterm{def:b3:probability:independence}{independientes},
el~\cref{exo:b3:probability:3}(a)), centradas y con $\V(Z_n) =
x_n^2$: la pregunta 13 se aplica cuando $\sum_nx_n^2 < \infty$ y da
convergencia c.s. En general, para cada $N$ la convergencia de
$\sum_nx_n\varepsilon_n$ no se ve afectada por los valores de
$\varepsilon_1, \dots, \varepsilon_N$: el suceso de convergencia está en
la $\sigma$-álgebra de cola de la sucesión
independiente
$(\varepsilon_n)$, de modo que la
\omterm{thm:b3:probability:zeroone}{ley cero--uno} de Kolmogorov
(el~\cref{thm:b3:probability:zeroone}) fuerza que su probabilidad valga
$0$ o $1$.

\textbf{15.} (a) Partiendo en el nivel $\theta\E Z$ y aplicando
Cauchy--Schwarz al trozo superior,
\[
\E Z = \E\bigl[Z\mathbf 1_{Z \leq \theta\E Z}\bigr] +
\E\bigl[Z\mathbf 1_{Z > \theta\E Z}\bigr]
\leq \theta\,\E Z + \sqrt{\E Z^2}\,
\sqrt{\P(Z > \theta\E Z)} ,
\]
de modo que $(1 - \theta)\E Z \leq \sqrt{\E Z^2\,\P(Z > \theta\E
Z)}$; elévese al cuadrado. (b) Desarróllese $T_n^4 =
\sum_{i,j,k,l}x_ix_jx_kx_l\,
\E[\varepsilon_i\varepsilon_j\varepsilon_k\varepsilon_l]$: la
\omterm{def:b3:probability:space}{esperanza} vale $1$ cuando los índices
se emparejan (los cuatro iguales, o dos pares distintos, esto último en
$3$ disposiciones) y $0$ en caso contrario (un signo desemparejado tiene
media nula y sale factor por
\omterm{def:b3:probability:independence}{independencia}). Por tanto,
\[
\E T_n^4 = \sum_kx_k^4 + 3\sum_{i\neq j}x_i^2x_j^2 =
3s_n^4 - 2\sum_kx_k^4 \leq 3s_n^4 .
\]
(c) Paley--Zygmund con $Z = T_n^2$, $\E Z = s_n^2$,
$\theta = \frac14$:
\[
\P\Bigl(\abs{T_n} > \frac{s_n}2\Bigr) = \P\Bigl(T_n^2 >
\frac{s_n^2}4\Bigr) \geq \Bigl(\frac34\Bigr)^2
\frac{s_n^4}{3s_n^4} = \frac3{16} .
\]
Si la serie convergiera con probabilidad positiva, convergería c.s.\
(pregunta 14), de modo que $\sup_n\abs{T_n} <
\infty$ c.s., y algún $M$ cumpliría
$\P(\sup_n\abs{T_n} > M) < \frac3{16}$; pero en cuanto $s_n
> 2M$, $\P(\abs{T_n} > M) \geq \P(\abs{T_n} > \frac{s_n}2)
\geq \frac3{16}$: contradicción. Así, la divergencia es casi segura y,
con la pregunta 14, la dicotomía es
\omterm{def:b3:complete:complete}{completa}.

\textbf{16.} Aquí $x_n = n^{-s}$ y $\sum_nn^{-2s} <
\infty$ exactamente cuando $s > \frac12$: por las preguntas 14--15,
$\sum_n\frac{\varepsilon_n}{n^s}$ converge c.s.\ si y solo si
$s > \frac12$ (para $s \leq \frac12$, divergencia c.s.). Para
$\frac12 < s \leq 1$ la convergencia nunca es absoluta. La comparación
es instructiva: unos signos perfectamente alternados cancelan con
intensidad $n^{-s}$ para todo $s
> 0$, mientras que unos signos aleatorios típicos cancelan solo con
intensidad raíz cuadrada --- el paseo aleatorio de la pregunta 21 crece
como $\sqrt n$, y la sumación de Abel convierte exactamente ese
crecimiento en convergencia de $\sum\varepsilon_nn^{-s}$ para
$s > \frac12$.

\textbf{17.} (a) $\cosh\lambda =
\sum_k\frac{\lambda^{2k}}{(2k)!}$ y $\eu^{\lambda^2/2} =
\sum_k\frac{\lambda^{2k}}{2^kk!}$; y $(2k)! \geq 2^kk!$ vale término a
término, porque $\frac{(2k)!}{k!} =
\prod_{i=1}^k(k + i) \geq \prod_{i=1}^k(2i) = 2^kk!$ (cada factor cumple
$k + i \geq 2i$ para $i \leq k$), de modo que, de hecho,
$(2k)! \geq 2^k(k!)^2 \geq 2^kk!$. (b) Obsérvese $a \leq 0 \leq b$ ($Z$
está centrada) y, por convexidad de $z \mapsto \eu^{\lambda z}$, para
$z \in \intcc ab$:
\[
\eu^{\lambda z} \leq \frac{b - z}{b - a}\,\eu^{\lambda a} +
\frac{z - a}{b - a}\,\eu^{\lambda b},
\qquad\text{luego}\qquad
\E\,\eu^{\lambda Z} \leq \frac{b\,\eu^{\lambda a} -
a\,\eu^{\lambda b}}{b - a}
= (1 - p)\eu^{-pt} + p\,\eu^{(1-p)t} = \eu^{\varphi(t)}
\]
con $p = \frac{-a}{b-a} \in \intcc01$, $t = \lambda(b -
a)$, $\varphi(t) = -pt + \log(1 - p + p\eu^t)$. Entonces
$\varphi(0) = 0$, $\varphi'(t) = -p + \frac{p\eu^t}{1 - p +
p\eu^t}$ se anula en $0$ y $\varphi''(t) = \rho(1 -
\rho) \leq \frac14$ para $\rho = \frac{p\eu^t}{1 - p +
p\eu^t} \in \intcc01$: Taylor de orden $2$ da
$\varphi(t) \leq \frac{t^2}8 = \frac{\lambda^2(b-a)^2}8$.

\textbf{18.} Para $\lambda > 0$, Markov aplicado a la variable positiva
$\eu^{\lambda(S_n - \E S_n)}$ (la~\cref{prop:b3:probability:markov}) y
la fórmula del producto para variables
\omterm{def:b3:probability:independence}{independientes} dan
\[
\P(S_n - \E S_n \geq t) \leq \eu^{-\lambda
t}\prod_{i=1}^n\E\,\eu^{\lambda(X_i - \E X_i)}
\leq \exp\Bigl(-\lambda t +
\frac{\lambda^2}8\sum_i(b_i - a_i)^2\Bigr),
\]
por la pregunta 17(b) aplicada a cada $X_i - \E X_i
\in \intcc{a_i - \E X_i}{b_i - \E X_i}$ centrada (misma anchura).
Minimizando el exponente en $\lambda = \frac{4t}{D}$, $D =
\sum_i(b_i - a_i)^2$, se obtiene $-\frac{2t^2}D$. La cola inferior se
sigue aplicando el resultado a $(-X_i)$.

\textbf{19.} Tómense $t = n\varepsilon$ y $D = n(b - a)^2$:
\[
\P\Bigl(\Bigl|\frac{S_n}n - m\Bigr| \geq \varepsilon\Bigr)
\leq 2\exp\Bigl(\frac{-2n^2\varepsilon^2}{n(b-a)^2}\Bigr)
= 2\exp\Bigl(\frac{-2n\varepsilon^2}{(b-a)^2}\Bigr),
\]
que es sumable en $n$ (una serie de tipo geométrico): Borel--Cantelli
(el~\cref{thm:b3:probability:borelcantelli}) da que c.s.\
$\abs{\frac{S_n}n - m} < \varepsilon$ a partir de cierto índice;
intersecando en $\varepsilon = \frac1p$ resulta $\frac{S_n}n \to m$ c.s.
Comparación: Etemadi solo pide $X_1 \in L^1$ e
\omterm{def:b3:probability:independence}{independencia} dos a dos, y no
entrega velocidad alguna; Hoeffding pide acotación e
\omterm{def:b3:probability:independence}{independencia} \omterm{def:b3:complete:complete}{completa}, y
entrega una garantía exponencial explícita en cada $n$ finito --- los
dos teoremas responden a preguntas distintas sobre el mismo límite.

\textbf{20.} Las $g(U_k)$ son i.i.d.\ con valores en $\intcc01$ y media
$\int g\,\dd\lambda_d$ (transferencia), de modo que la pregunta 18 con
$b_i - a_i = 1$, $t = n\varepsilon$ da la cota bilátera
$2\eu^{-2n\varepsilon^2} \leq \delta$ en cuanto
$\eu^{2n\varepsilon^2} \geq \frac2\delta$, es decir,
$n \geq \frac{\log(2/\delta)}{2\varepsilon^2}$. Para
$\varepsilon = \delta = 10^{-2}$:
\[
n \geq \frac{\log 200}{2\cdot10^{-4}} =
\frac{5.2983\ldots}{0.0002} \approx 26\,492 :
\]
unas $26\,500$ muestras garantizan una precisión $1\%$ con confianza
$99\%$ --- en toda dimensión $d$ y para todo integrando
\omterm{def:b3:lebesgue:measurable}{medible} con valores en $\intcc01$.
La \omterm{def:b3:probability:space}{ley} fuerte de la pregunta 11
prometía convergencia sin ninguna garantía a $n$ finito; una malla
determinista con $k$ puntos por eje cuesta $k^d$ evaluaciones,
exponencial en $d$. La concentración es lo que hace de
\omterm{ex:b3:probability:sllnapps}{Monte Carlo} un \emph{método} y no
una \omterm{def:b3:probability:space}{esperanza}.

\textbf{21.} \omterm{def:b3:probability:independence}{Independencia} y
fórmula del producto:
$\E\,\eu^{\lambda S_n} = (\E\,\eu^{\lambda\varepsilon_1})^n
= (\cosh\lambda)^n \leq \eu^{n\lambda^2/2}$ por la pregunta 17(a).
Markov sobre $\eu^{\lambda S_n}$:
\[
\P(S_n \geq x) \leq \eu^{-\lambda x + n\lambda^2/2}
= \eu^{-x^2/(2n)}
\qquad\text{en el óptimo } \lambda = \frac xn,
\]
y la cota simétrica para $-S_n$ (misma
\omterm{def:b3:probability:space}{ley}) duplica la constante para
$\abs{S_n}$.

\textbf{22.} Fíjese $\eta > 0$ y póngase $x_n = (1 +
\eta)\sqrt{2n\log n}$ para $n \geq 2$:
\[
\P(\abs{S_n} \geq x_n) \leq 2\exp\bigl(-(1 +
\eta)^2\log n\bigr) = \frac{2}{n^{(1+\eta)^2}},
\]
sumable, pues $(1 + \eta)^2 > 1$. Borel--Cantelli: c.s.\
$\abs{S_n} < (1 + \eta)\sqrt{2n\log n}$ para todo $n$ grande, de modo
que $\limsup_n\frac{\abs{S_n}}{\sqrt{2n\log n}} \leq 1 +
\eta$ c.s.; intersecando los sucesos c.s.\ para $\eta =
\frac1p$, $p \in \N^*$, se obtiene la afirmación. El paseo de tamaño $n$
tiene amplitud típica $\sqrt n$ (su varianza), y hasta sus peores
excursiones superan esa escala a lo sumo en
$\sqrt{2\log n}$.

\textbf{23.} Con $n_j = 2^j$ y $x = (1 +
\eta)\sqrt{2n_j\log\log n_j}$ (definido para $j \geq 2$), la pregunta 21
da
\[
\P\bigl(S_{n_j} \geq x\bigr) \leq \exp\bigl(-(1 +
\eta)^2\log\log n_j\bigr) = (j\log 2)^{-(1+\eta)^2},
\]
sumable en $j$, pues $(1 + \eta)^2 > 1$: Borel--Cantelli y
$\eta = \frac1p$ dan $\limsup_jS_{n_j}/\sqrt{2n_j
\log\log n_j} \leq 1$ c.s. Lo que falta para la mitad superior \omterm{def:b3:complete:complete}{completa}
es el puente entre los instantes de control: hay que demostrar que
$\max_{n_j\leq n\leq n_{j+1}}S_n$ supera
$(1+\eta)\sqrt{2n_j\log\log n_j}$ solo un número finito de veces, lo que
exige una desigualdad maximal con colas \emph{gaussianas} (la
desigualdad de reflexión de Lévy o la de Ottaviani, no demostradas
aquí). La pregunta 12 es cuantitativamente demasiado débil: acota la
probabilidad por
\[
\frac{n_j}{(1+\eta)^2\,2n_j\log\log n_j}
= \frac{1}{2(1+\eta)^2\log(j\log2)},
\]
que tiende a $0$ pero \emph{no es sumable} en $j$: Borel--Cantelli no
puede concluir. La mitad inferior de la
\omterm{def:b3:probability:space}{ley} del logaritmo iterado aplica el
segundo lema de Borel--Cantelli a los incrementos
\omterm{def:b3:probability:independence}{independientes}
$S_{n_{j+1}} - S_{n_j}$, usando cotas inferiores ajustadas para colas de
tipo gaussiano. Ambos refinamientos son probabilidad genuina de tercer
año, un curso más allá; lo que este problema entrega sin ayuda es la
escala exacta del logaritmo iterado a lo largo de tiempos geométricos.

\textbf{24.} Cada $\hat p_i$ es una media de $n$ variables indicadoras
i.i.d.\ con valores en $\intcc01$ y media $\P(A_i)$: Hoeffding da
$\P(\abs{\hat p_i - \P(A_i)} >
\varepsilon) \leq 2\eu^{-2n\varepsilon^2}$. La cota de la unión
multiplica por $N$. Resolviendo $2N\eu^{-2n\varepsilon^2} \leq
\delta$: $n \geq \frac{\ln(2N/\delta)}{2\varepsilon^2}$. Numéricamente:
$\ln\frac{2\cdot10^6}{0.05} =
\ln(4\cdot10^7) \approx 17.5$, de modo que $n \geq
\frac{17.5}{2\cdot10^{-4}} \approx 87\,600$: estimar \emph{una}
probabilidad con precisión $\pm1\%$ cuesta unas $18\,500$ muestras
($\ln(2/\delta)/2\varepsilon^2$), y \emph{un millón} de probabilidades
solo $\approx 4.7$ veces más --- la uniformidad cuesta $\ln N$, no $N$:
la observación que hace estadísticamente posible la minimización del
riesgo empírico y, con ella, el aprendizaje automático.

\textbf{25.} Las variables $X_n = \frac{\varepsilon_n}
{n^\alpha}$ son
\omterm{def:b3:probability:independence}{independientes}, centradas y
acotadas, con $\sum_n\V(X_n) = \sum_nn^{-2\alpha}$. Si $\alpha >
\frac12$: la serie de varianzas converge, y el teorema de la serie única
(Parte VI) da la convergencia c.s.\ de $\sum X_n$. Si
$\alpha \leq \frac12$: la serie de varianzas diverge, y la mitad
recíproca (el argumento de Paley--Zygmund de la Parte VI, aplicable
porque los sumandos están acotados por $1$) da la divergencia c.s. La
convergencia absoluta pide $\sum n^{-\alpha} < \infty$: $\alpha > 1$. En
$\intoc{\frac12}1$, la serie converge c.s.\ aunque
$\sum\abs{X_n} = \infty$ con seguridad: los signos conspiran para
cancelarse, con probabilidad uno --- convergencia por cancelación,
invisible para cualquier criterio absoluto y, por la
\omterm{thm:b3:probability:zeroone}{ley cero--uno}, con un veredicto
determinista de todos modos.
\end{solution}
